\documentclass[aspectratio=169]{beamer}
\input{header}
\coursecode{JKSSB}

\title{PowerPoint Basics}
\subtitle{Lesson 31 --- Slides, Layouts, Masters and Views}
\author{JKSSB Computer Awareness}
\date{\today}

\begin{document}
\frame{\titlepage}

\begin{frame}{What MS PowerPoint Is}
  \begin{itemize}
    \item \key{MS PowerPoint} is the presentation program of the MS Office suite.
    \item A \key{presentation} is a file made of one or more pages; each page
      is called a \key{slide}.
    \item Slides are shown one after another as a \key{slide show} on screen
      or on a projector.
    \item The exam name is \key{Powerpoint Presentation}, and the usual short
      form is \key{PPT}.
  \end{itemize}
  \begin{block}{The one idea}
    PowerPoint = slides shown in sequence; one slide is one screen of the show.
  \end{block}
\end{frame}

\begin{frame}{Full Forms and File Extensions}
  \begin{itemize}
    \item \key{PPT} --- PowerPoint Presentation; \key{PPTX} --- PowerPoint XML
      Presentation (default since Office 2007).
    \item \key{PPSX} --- PowerPoint Slide Show; \key{POTX} --- PowerPoint
      Template; \key{PPTM} --- macro-enabled presentation.
    \item Default working file: \key{.pptx}; template file: \key{.potx};
      show file: \key{.ppsx}.
    \item Presentations can also be saved or exported as \key{.pdf}.
  \end{itemize}
  \begin{alertblock}{Trap alert}
    \key{.pptx} is the editable presentation; \key{.ppsx} opens directly as a
    slide show; \key{.potx} is the template.
  \end{alertblock}
\end{frame}

\begin{frame}{Slide and Placeholder}
  \begin{itemize}
    \item A \key{slide} is a single page of a presentation; it can hold text,
      pictures, charts, tables, audio, and video.
    \item A \key{placeholder} is the dotted-line box on a slide that holds
      such content (title, text, picture, chart, table).
    \item Click inside a placeholder and type, or use its icon to insert a
      picture, chart, or table.
    \item Placeholders come from the slide's \key{layout}; resizing or moving
      the placeholder changes only that slide.
  \end{itemize}
  \begin{block}{Keep the pair straight}
    The slide is the page; the placeholder is the dotted box that holds content
    on that page.
  \end{block}
\end{frame}

\begin{frame}{Slide Layouts}
  \begin{itemize}
    \item A \key{layout} is the arrangement of placeholders on a slide.
    \item Common layouts: Title Slide, Title and Content, Section Header, Two
      Content, Comparison, Blank.
    \item Apply or change a layout from Home tab $\rightarrow$ \key{Layout};
      adding a new slide uses Home tab $\rightarrow$ \key{New Slide}.
    \item The \key{Blank} layout has no placeholders; content must be inserted
      manually.
  \end{itemize}
  \begin{block}{Office desktop example}
    A class report uses Title Slide once, then Title and Content for every
    following point --- same layout family, different content.
  \end{block}
\end{frame}

\begin{frame}{Themes vs Templates}
  \begin{columns}[T]
    \column{0.46\textwidth}
      \textbf{Theme}
      \begin{itemize}
        \item A coordinated set of colours, fonts, and effects.
        \item Applied from the Design tab; one click restyles all slides.
      \end{itemize}
    \column{0.46\textwidth}
      \textbf{Template}
      \begin{itemize}
        \item A ready-made pattern file (\key{.potx}) with a theme plus
          layouts and sample content.
        \item Starting from a template reuses its design for a new file.
      \end{itemize}
  \end{columns}
  \vspace{0.3cm}
  \muted{A theme restyles; a template is a reusable starter file.}
\end{frame}

\begin{frame}{Slide Master and the Three Masters}
  \begin{itemize}
    \item The \key{Slide Master} is the top-level slide that controls the
      design of all layouts and slides linked to it.
    \item Change the font, logo, or footer on the Slide Master and every
      linked slide follows.
    \item PowerPoint has three masters: \key{Slide Master}, \key{Handout
      Master}, and \key{Notes Master}.
    \item Open it from View tab $\rightarrow$ \key{Slide Master}; close with
      Close Master View.
  \end{itemize}
  \begin{alertblock}{Trap alert}
    Edit once on the Slide Master instead of fixing the same logo on forty
    separate slides.
  \end{alertblock}
\end{frame}

\begin{frame}{Views at a Glance}
  \begin{center}
    \begin{tabular}{p{0.22\textwidth}p{0.62\textwidth}}
      \toprule
      \textbf{View} & \textbf{Shows} \\
      \midrule
      Normal & Default editing view: slide pane + thumbnails + notes pane \\
      Outline View & Text outline of all slides \\
      Slide Sorter & Thumbnails of all slides; reorder and organise \\
      Reading View & Full-screen-like reading without editing tools \\
      Slide Show & Real full-screen show from a chosen slide \\
      Notes Page & Slide above, speaker notes below \\
      \bottomrule
    \end{tabular}
  \end{center}
  \vspace{0.25cm}
  \muted{Normal is the default view; Slide Sorter manages order.}
\end{frame}

\begin{frame}{Normal View: the Default}
  \begin{itemize}
    \item \key{Normal view} is the default editing view of PowerPoint.
    \item Three working areas: the large \key{slide pane}, the
      \key{slides/outline pane} (thumbnails on the left), and the
      \key{notes pane} below.
    \item All typing, formatting, and inserting happen here.
    \item Switch views from the View tab or the buttons on the status bar.
  \end{itemize}
  \begin{block}{The one idea}
    If the question says ``default editing view'', the answer is Normal view.
  \end{block}
\end{frame}

\begin{frame}{Slide Sorter and Reading Views}
  \begin{itemize}
    \item \key{Slide Sorter} shows thumbnails of every slide together, so
      slides can be reordered by drag and drop.
    \item It is the best view for deleting, duplicating, or checking the whole
      order at a glance.
    \item \key{Reading View} shows the show almost full-screen for reading,
      with simple navigation, but without the editing tools.
    \item Reading View suits review on the author's own screen; Slide Show
      suits projection to an audience.
  \end{itemize}
\end{frame}

\begin{frame}{Notes Page and Speaker Notes}
  \begin{itemize}
    \item \key{Notes} (speaker notes) are reminders typed in the notes pane
      below a slide; the audience never sees them during the show.
    \item \key{Notes Page view} shows one slide above its notes text, one page
      per slide.
    \item Notes print through File $\rightarrow$ Print $\rightarrow$
      \key{Notes Pages}.
    \item During the show, notes are visible only in \key{Presenter View} on
      the speaker's screen.
  \end{itemize}
  \begin{block}{Keep the pair straight}
    Notes help the speaker; handouts help the audience.
  \end{block}
\end{frame}

\begin{frame}{Handouts and Printing}
  \begin{itemize}
    \item \key{Handouts} are printed sheets with several slides per page (1, 2,
      3, 4, 6, or 9 slides).
    \item The common exam choice is \key{6 slides per page}; 3-per-page adds
      ruled lines for the audience to write on.
    \item Handout design is controlled by the \key{Handout Master}.
    \item Print path: File $\rightarrow$ Print; options include Full Page
      Slides, Handouts, Notes Pages, and Outline View.
  \end{itemize}
\end{frame}

\begin{frame}{Running the Slide Show}
  \begin{itemize}
    \item \key{F5} starts the show from the beginning; \key{Shift+F5} starts
      from the current slide.
    \item Move forward with N, Space, Right arrow, or Enter; go back with P or
      Left arrow; end with \key{Esc}.
    \item Press \key{B} for a black screen and \key{W} for a white screen
      during the show.
    \item The Slide Show tab sets timing, rehearsal, and monitor options.
  \end{itemize}
  \begin{alertblock}{Trap alert}
    F5 = from the beginning; Shift+F5 = from the current slide. (WO 2026 Q80)
  \end{alertblock}
\end{frame}

\begin{frame}{Transitions vs Animations}
  \begin{columns}[T]
    \column{0.46\textwidth}
      \textbf{Transition}
      \begin{itemize}
        \item Movement \key{between} slides (Fade, Push, Wipe, Morph).
        \item Set on the Transitions tab; one per slide.
      \end{itemize}
    \column{0.46\textwidth}
      \textbf{Animation}
      \begin{itemize}
        \item Movement of \key{objects on} a slide: Entrance, Emphasis, Exit,
          Motion Paths.
        \item Set on the Animations tab; managed in the Animation Pane.
      \end{itemize}
  \end{columns}
  \vspace{0.3cm}
  \begin{alertblock}{Trap alert}
    Transition moves the slide; animation moves the object on the slide.
  \end{alertblock}
\end{frame}

\begin{frame}{Trap Alert: PowerPoint Facts to Keep Straight}
  \begin{itemize}
    \item \key{.pptx} = editable presentation; \key{.ppsx} = show file;
      \key{.potx} = template.
    \item Slide Master controls all linked slides; Handout and Notes Masters
      control print layouts.
    \item \key{Normal} is the default view; \key{Slide Sorter} reorders.
    \item Notes are for the \key{speaker}; handouts are for the
      \key{audience}.
    \item F5 starts from the \key{beginning}; Shift+F5 from the
      \key{current} slide.
  \end{itemize}
\end{frame}

\begin{frame}{Lesson 31: Recall}
  \begin{itemize}
    \item Presentation = ordered slides; each slide holds placeholders placed
      by its layout.
    \item Theme = colours, fonts, effects; template = reusable starter file
      (.potx); Slide Master controls all linked slides.
    \item Views: Normal (default editing), Slide Sorter (reorder), Reading
      (review), Notes Page (speaker notes), Slide Show (present).
    \item Notes help the speaker; handouts (1--9 per page) help the audience;
      print via File $\rightarrow$ Print.
    \item F5 from the beginning, Shift+F5 from the current slide; transitions
      move between slides, animations move objects on a slide.
  \end{itemize}
\end{frame}
\end{document}
